\section{Love and Memory Made Whole}

Our desire for heaven often begins with particular people. The death of a parent, a spouse, a child, or a friend makes eternal life cease to be merely a distant subject. We want to know whether the love shared in time will be lost. In the communion of saints and the resurrection, the same persons are saved and their love is perfected in God. The fulfillment he promises exceeds our ability to preserve the familiar terms on which we first learned to love.

Christ speaks plainly about one such transformation: the risen neither marry nor are given in marriage (Mt 22:30). The Sadducees' question imagines resurrection as the continuation of present marital claims under impossible circumstances. Jesus directs them to God's power and to the life he gives. Marriage is an earthly vocation with its own holiness, fidelity, and sacramental meaning; it is not an institution perpetuated in the resurrection. The love formed through it is entrusted to the fullness of charity. A husband and wife can hope to love one another perfectly in God without imagining their earthly marriage continuing in another form.

Trusting this promise requires us to surrender some present expectations about affection. Our present love is intertwined with exclusivity, need, habit, and sometimes possession. Even a generous love can be afraid that the beloved's wider happiness will reduce its own place. Heavenly charity heals that fear. The good of another person is loved in God, whose fullness does not compete with our love but makes it truthful. We do not know the exact texture of every future relationship. We do know that perfected charity will contain no jealousy, no manipulation, and no fear that God's gift to someone else deprives us of the good he wills for us.

Cyprian's hope of reunion has strength within this larger charity. He names the embrace of dear ones and the gladness shared between them, then places that reunion among the apostles, martyrs, and all the friends of God (\emph{On Mortality} 26). Cyprian's exhortation concerns the reunion of the saved; the final state of a particular departed person remains in God's judgment. We may pray and hope for the salvation of those we love. Their presence in God would perfect the love we have begun, placing it within a communion greater than we can now receive.

The promise also concerns memory. Much of our present self-understanding is bound to what has hurt us or to what we have done wrong. We can fear that happiness would require the erasure of our history, while remembering seems inseparable from renewed grief or shame. Augustine considers this tension in \emph{City of God} XXII.30. He distinguishes intellectual knowledge of an evil from continued experience of its suffering. The blessed know their deliverance and give thanks to the Redeemer, while the misery from which they were delivered no longer afflicts them. Their gratitude is truthful, and their healing is complete.

Augustine illustrates the distinction through the knowledge of a physician who understands diseases he has not personally suffered. His purpose is to show that knowledge need not reproduce affliction within the knower. This is his theological explanation of blessed memory; revelation does not give us a detailed account of every remembered event. The divine promise is that tears, mourning, and death will pass away (Rev 21:4). We can therefore entrust our history to God without deciding beforehand which particular mode of memory must constitute its healing. The Savior who knows our lives more truly than we do can bring them to peace.

The hope of healed memory has consequences now. We can repent without treating shame as an eternal identity. We can tell the truth about an injury without making bitterness the proof that we have taken it seriously. Forgiveness does not require the denial of justice, and grief does not have to be forced into an appearance of ease. But neither our guilt nor our wounds can define the limits of Christ's saving power. To place the past before him is to consent to a future in which it is known within mercy and truth. We ask to be healed in the truth of our lives, trusting that the misery we have known need not remain our condition.

